\section{Regional Climate and the General Circulation}
\subsection{The Atmosphere}
An important aspect of the Arctic climate is the connection to the atmospheric circulation through the large-scale circulation patterns known as the Hadley, Ferrel, and Polar cells. The Polar cell is thus coupled with the general circulation and lifts air masses from the surface to the upper troposphere, moving them toward the North Pole, where they subside, creating a high-pressure system (as seen in \autoref{fig:SeaLevelPressure_mean}) \citep{hartmann2016}. This creates a dry desert-like climate and easterly winds, which also help drive the water southward through Ekman transport. The strength of the pressure difference between the Arctic high pressure and the mid-latitudes can be measured using the Arctic Oscillation (AO) index, which is directly linked to the North Atlantic Oscillation (NAO) \citep{serreze2014}. In \autoref{fig:SeaLevelPressure_mean_DJF} and \autoref{fig:SeaLevelPressure_mean_JJA} one can see that the AO index is positive during the winter, while weaker during the summer. This is because the AO index is directly linked to the polar cell, which is linked to the temperature gradient just like the Hadley cell, thus creating stronger AO during the winter. 

\subsection{The Ocean}
The Arctic Ocean is connected to the Atlantic Ocean through the Fram Strait but is otherwise almost disconnected from the world's oceans. Thus, most oceanic water entering and exiting the Arctic Ocean passes through the Fram Strait (as seen in \autoref{fig:SeaSurfaceTemperature_annual}) \citep{serreze2014}. The incoming water is the Norwegian Current, which is a branch of the Gulf Stream and thus warm and saline. When this water reaches higher latitudes, it is much warmer than the atmosphere, leading to a large heat flux from the ocean to the atmosphere. Furthermore, the water becomes less saline as it mixes with freshwater from river discharge. As some water freezes, brine rejection leads to an increase in salinity and thus an increase in potential density, causing the water to sink. Interestingly, this creates an increase in temperature with depth for the initial hundred meters, which is the opposite of other oceans \citep{serreze2014}.

Thus, the water that entered the Arctic Ocean through the Fram Strait eventually leaves the Arctic Ocean through the Fram Strait, but now as cold and less saline bottom water known as North Atlantic Deep Water (NADW). These processes are an important part of the global meridional overturning circulation. Additionally, as the region experiences large river discharge, more ocean water must leave the Arctic Ocean to balance the water budget. This, together surface winds, creates a southward transport of cold, fresh Arctic surface water through the East Greenland Current, which continues southward and contributes to the Labrador Current \citep{serreze2014}. Furthermore, sea ice is also exported to the North Atlantic through the Beaufort Gyre, a slow-moving system that circulates sea ice in a clockwise direction around the North Pole \citep{serreze2014}. This process creates layers of old sea ice along the Arctic Archipelago and also exports sea ice through the Fram Strait.